\graphicspath{{chapters/contents/sessio_01/images/}}

\section{Sessió 02 (21/09/26). Quadern de laboratori: absorció de la llum solar segons el color}

\begin{tcolorbox}[
    colback=white,
    colframe=black!40,
    boxrule=0.6pt,
    arc=2mm,
    left=3mm,
    right=3mm,
    top=2mm,
    bottom=2mm
]
\textbf{Nota de seguretat.} No s’ha de mirar mai, sota cap concepte, directament al Sol, a un làser, a una font de llum ultraviolada o a qualsevol font de llum intensa. Els danys a l’ull són irreversibles en la majoria dels casos i poden arribar fins a la pèrdua permanent de visió.
\end{tcolorbox}



En aquest experiment exposarem al Sol pedres idèntiques en forma i mida, pintades de blanc, gris i negre, i en mesurarem la temperatura després d’un temps d’exposició. També deixarem tres pedres de cada color guardades dins d’una caixa, sense exposició directa al Sol, com a grup de control.

Anota els resultats i, abans de llegir l’explicació del final d’aquesta fitxa, intentarem descobrir entre tots i totes què ha passat i per què.

\subsection{Material}

\begin{itemize}

    \item 1 pedra blanca, 1 pedra grisa i 1 pedra negra per al grup de control (dins de la caixa).

    \item 1 pedra blanca, 1 pedra grisa i 1 pedra negra per al grup exposat al Sol.

    \item Pantalla transparent paravent (per exemple, de metacrilat o una caixa de plàstic transparent), per evitar que el vent refredi les pedres per convecció sense bloquejar la llum.

    \item Suports que mantinguin les pedres separades del terra, per evitar que guanyin o perdin calor per conducció en tocar-lo.

    \item Termòmetre làser (termòmetre infraroig sense contacte).

    \item Full o taula per al registre de dades.

\end{itemize}

\subsection{Procediment}

\textbf{Pas 1.} Prepara una pedra de cada color dins de la caixa, protegides de la llum solar directa. Aquest és el grup de control.

\textbf{Pas 2.} Prepara tres pedres més, una de cada color en un lloc exposat al Sol. És molt important controlar tot allò que pugui alterar el resultat: el vent pot refredar les pedres, de manera que cal buscar un lloc protegit o utilitzar la pantalla transparent paravent; el terra pot transmetre o absorbir calor per contacte, així que cal utilitzar els suports per separar-ne les pedres.

\textit{Aïllar bé les condicions i controlar les variables que no volem que influeixin és una part essencial de qualsevol experiment.}

\textbf{Pas 3.} Deixa passar 30 minuts perquè les pedres exposades al Sol arribin a l’equilibri tèrmic i s’hagin escalfat tant com sigui possible.

\textbf{Pas 4.} Mesura la temperatura de les tres pedres de la caixa (grup de control) i anota-la a la taula 1. Segons les indicacions del fabricant, el termòmetre ha d’estar a uns 25 cm. Assegura’t de tenir ben interioritzada aquesta distància.

\textbf{Pas 5.} Mesura la temperatura de les tres pedres exposades al Sol. Durant la mesura, posa temporalment cada pedra a l’ombra perquè el termòmetre infraroig pugui mesurar-ne correctament la temperatura sense que la radiació solar interfereixi en la lectura. Repeteix la mesura tres vegades per pedra i anota les dades a la taula 2.

\textbf{Pas 6.} Calcula la mitjana de les tres mesures de cada pedra exposada al Sol.

\subsection{Resultats}


% Tabla 1
\renewcommand{\arraystretch}{2.4}

\begin{table}[h]
\centering
\caption{Temperatura del grup de control (pedres dins de la caixa)}

\begin{tabular}{|>{\centering\arraybackslash}m{1.5cm}|>{\centering\arraybackslash}m{3.3cm}|}
\hline
\textbf{Color} & \textbf{Temperatura ($^\circ$C)} \\
\hline
Blanc & \\
\hline
Gris & \\
\hline
Negre & \\
\hline
\end{tabular}

\end{table}

\renewcommand{\arraystretch}{1}

% Tabla 2
\renewcommand{\arraystretch}{2.4}
\setlength{\tabcolsep}{4pt}

\begin{table}[h]
\centering
\caption{Temperatura de les pedres exposades al Sol}

\begin{tabular}{|
>{\centering\arraybackslash}m{1.2cm}|
>{\centering\arraybackslash}m{1.5cm}|
>{\centering\arraybackslash}m{1.5cm}|
>{\centering\arraybackslash}m{1.5cm}|
>{\centering\arraybackslash}m{1.5cm}|
}
\hline
\textbf{Color} &
\shortstack{\textbf{T 1}\\\textbf{($^\circ$C)}} &
\shortstack{\textbf{T 2}\\\textbf{($^\circ$C)}} &
\shortstack{\textbf{T 3}\\\textbf{($^\circ$C)}} &
\shortstack{\small{\textbf{Mitjana}}\\\textbf{($^\circ$C)}} \\
\hline
Blanc & & & & \\
\hline
Gris & & & & \\
\hline
Negre & & & & \\
\hline
\end{tabular}

\end{table}

\setlength{\tabcolsep}{6pt}
\renewcommand{\arraystretch}{1}

\subsection{Discussió}

\begin{itemize}
    \item Quina de les pedres exposades al Sol s’ha escalfat més? I quina menys?
    \item Com es comparen aquestes temperatures amb les del grup de control?
    \item Se t’acut alguna explicació de per què ha passat això?
    \item Per què calia controlar el vent i el contacte amb el terra? Què hauria passat si no s’hagués fet?
    \item L’experiment ha sortit com esperaves? Com el podries millorar o repetir amb més precisió?
\end{itemize}

\subsection{Explicació}

La llum del Sol està formada per diferents tipus de radiació. Una part és la llum visible, és a dir, la que podem veure en forma de colors. Altres tipus de radiació, com la infraroja o la ultraviolada, no els podem veure, però també transporten energia.

Quan un objecte absorbeix aquesta energia, s’escalfa i la seva temperatura augmenta. El color d’un objecte depèn de la llum que reflecteix i de la que absorbeix. Un objecte blanc reflecteix gran part de la llum que rep, mentre que un objecte negre n’absorbeix molta més. Per això, normalment, un objecte negre s’escalfa més que un de blanc quan està exposat al Sol.

Per això, les pedres negres s’haurien d’haver escalfat més que les grises, i aquestes més que les blanques, mentre que les pedres de la caixa, com que no han rebut llum solar directa, s’haurien d’haver mantingut totes a una temperatura molt semblant entre elles.

\subsection{Activitats}

\subsubsection{Nivell I. Podrien repetir el nostre experiment?}

Un altre grup no ha vist com hem fet l’experiment. La seva feina serà intentar repetir-lo utilitzant només el teu esquema.

Dibuixa el muntatge experimental de manera que puguin entendre què hem fet sense haver de preguntar-te res.

Al dibuix han d’aparèixer:

\begin{itemize}
    \item les pedres exposades al Sol i les del grup de control,
    \item les temperatures que hem mesurat,
    \item la pantalla paravent i els suports,
    \item qualsevol indicació que consideris necessària perquè l’experiment es pugui repetir correctament.
\end{itemize}

Calcula també, per a cada color, quant s’ha escalfat la pedra respecte del seu control:

\[
\Delta T = T_{\text{Sol}} - T_{\text{control}}
\]

Anota aquests valors directament sobre el dibuix.

\textit{En ciència és fonamental comunicar les idees i els resultats amb claredat perquè tota la comunitat els pugui entendre, revisar i comprovar.}

\subsubsection{Nivell II. Quant costa el color?}

Un dipòsit industrial de 100\,000 litres conté nitrogen gasós, que s’ha de mantenir a una temperatura màxima de $20\,^\circ$C.

En un dia assolellat, el dipòsit s’escalfa per l’acció del Sol. Suposa que la temperatura que assoleix depèn del color amb què està pintat i utilitza, per a cada color, les temperatures que has obtingut en l’experiment amb les pedres.

Refredar el gas té un cost diari de 0,02~€/litre per cada grau que cal abaixar la temperatura.

\begin{enumerate}[label=\alph*)]
    \item Calcula quant costaria refredar el dipòsit fins als $20\,^\circ$C durant aquell dia si fos blanc, gris o negre.
    
    \item Calcula quants diners s’estalviaria l’empresa en un dia assolellat si pintés el dipòsit de blanc en lloc de negre.
\end{enumerate}

